\documentclass[11pt,a4paper]{article}

\usepackage[utf8]{inputenc}
\usepackage{cmap}
\usepackage[T1]{fontenc}
\usepackage[ngerman]{babel}
\usepackage{lmodern}
\usepackage{microtype}
\usepackage{amsmath,amssymb,amsthm}
\usepackage{mathtools}
\usepackage{enumitem}
\usepackage{hyperref}
\usepackage[nameinlink,capitalize]{cleveref}
\usepackage[margin=1in]{geometry}
\usepackage{booktabs}
\usepackage{array}
\usepackage{tabularx}
\usepackage{graphicx}
\usepackage{tikz}
\usetikzlibrary{arrows.meta,positioning,shapes.geometric,fit}

\input{glyphtounicode}
\pdfgentounicode=1

\hypersetup{
  unicode=true,
  pdflang={de-DE},
  colorlinks=true,
  linkcolor=blue,
  citecolor=blue,
  urlcolor=blue,
  pdftitle={Die Bealsche Vermutung via Height Dominance mit einem Potenzrest-Sieb-Begleiter},
  pdfauthor={Lukas Geiger},
  pdfsubject={Bedingte Beweisarchitektur zur Bealschen Vermutung mit Potenzrest-Sieb-Begleiter},
  pdfkeywords={Bealsche Vermutung, Höhendominanz, Height Dominance, Radikal, Potenzgleichungen, primitive ganzzahlige Tripel, Mikrocluster-Architektur, Potenzrest-Sieb, Large Sieve, Charaktersummen}
}

\DeclareMathOperator{\rad}{rad}
\DeclareMathOperator{\supp}{supp}
\emergencystretch=3em

\newtheorem{theorem}{Theorem}[section]
\newtheorem{lemma}[theorem]{Lemma}
\newtheorem{proposition}[theorem]{Proposition}
\newtheorem{conjecture}[theorem]{Vermutung}
\newtheorem{definition}[theorem]{Definition}
\newtheorem{remark}[theorem]{Bemerkung}

\crefname{theorem}{Theorem}{Theoreme}
\crefname{lemma}{Lemma}{Lemmata}
\crefname{proposition}{Proposition}{Propositionen}
\crefname{conjecture}{Vermutung}{Vermutungen}
\crefname{definition}{Definition}{Definitionen}
\crefname{remark}{Bemerkung}{Bemerkungen}

\title{Die Bealsche Vermutung via Height Dominance:\\
Eine Zookeeper-artige Beweisarchitektur\\
mit einem Potenzrest-Sieb-Begleiter}
\author{Lukas Geiger}
\date{Oktober 2026 --- Version 1.1 (Korrekturfassung)}

\begin{document}
\maketitle

\begin{abstract}
\enlargethispage{5\baselineskip}
\setlength{\parskip}{0.5ex}
Dieser Beitrag ist eine bedingte Architektur zur Bealschen Vermutung, kein
Beweis. Wir untersuchen die Vermutung entlang zweier verwandter Linien,
die einen gemeinsamen arithmetischen Input teilen.

\emph{Teil A (deterministische Architektur).}
Wir formulieren eine direkte Height-Dominance-Architektur (Höhendominanz-Architektur), die bewusst dem strukturellen Muster der Zookeeper-artigen Reduktion zur Riemannschen Vermutung nachempfunden ist: eine intrinsische Cluster-Definition, ein Cancellation-Koordinatensystem und ein Drei-Lemma-Endspiel, zentriert um ein zwingendes Komplement HD3'. Die primitive Potenzstruktur liefert die elementare Kompressionsschranke $\rad(ABC) < H^{1/x+1/y+1/z}$ ($H=C^z$),
und das noch fehlende zwingende Verifier-Komplement (Coercive Verifier-Complement) HD3' würde einen gemeinsamen Primteiler von $A,B,C$ erzwingen. HD3' wird als
\emph{Beal-starke lokal-globale Vermutung} formuliert, nicht als
schwächeres technisches Residuum: Sie konzentriert im Wesentlichen den
gesamten arithmetischen Gehalt der Bealschen Vermutung in eine einzige
scharf formulierte Aussage, sobald HD1 (skalare Höhen-Aufhebung) und
HD-PC (Potenz-Kompression) jedes hypothetische Gegenbeispiel in einem
kontrollierten Höhenregime platziert haben. Der bedingte Abschluss ist
daher eine \emph{Reduktionsarchitektur} und kein eigenständiger
Fortschrittsschritt in Richtung Beal. Das Motiv der entgegengesetzten
Radikal-Ungleichung erscheint nur als Motivation; der eigentliche
Abschluss-Mechanismus ist der lokal-globale Übergang HD3'. Das offene
Problem besteht darin, HD3' mit arithmetischen Mitteln zu etablieren.

Die Guardrail-Audits vom Juni 2026 werten diesen Status nicht auf: Es
wurde kein natürlicher Q97-plus-Survivor und kein direkter
S-Unit-/Frey-Anker gefunden; synthetische lokale Pässe kalibrieren nur
die Negativkontrollen des Root-Coherence-Gates.

\emph{Teil B (asymptotischer Sieb-Begleiter).}
Für positive primitive Lösungen mit festem $K\ge4$ und $3\le x,y,z\le K$,
ohne $(3,3,3)$, ergibt die korrigierte positive Zeilenrekonstruktion auf
klassischen Jacobi-/Weil- und Großsieb-Eingaben
$T_K(H)\ll_{K,\varepsilon}H^{2/3}/\Sigma(Q;K)$ für hinreichend große $H$
und $(\log H)^{2+\varepsilon}\le Q\le H^{1/(2K)}$.
Der untere Cut ist hinreichend, nicht notwendig. Die ganzzahlige Nullzeile
wird ausgeschlossen; Nullrestzeilen, leere Klassen und ein verschwindendes
Ausschlussbudget werden ausdrücklich behandelt. Keine explizite Schwelle
$H_0(K,\varepsilon)$ und kein effektives Konstantenpaket sind hier zertifiziert.
Das formale HD2'-Entropiebudget folgt asymptotisch aus Dichten für feste
Exponenten und Primzahlen in festen arithmetischen Progressionen; es ist
von den Basen unabhängig und beweist keine Wurzelkohärenz.
Die exponentielle Schranke der kumulativen positiven Lösungszählung würde
unter diesen Entropie-Eingaben $T_K\equiv0$ erzwingen; für alle festen $K$
ist sie bereits Beal-stark und keine schwache Zwischenschranke.

Version 1.1 berichtigt Genus- und Regime-Klassifikation, Nullzeilenrichtung
und Siebdomänen, skalare Monodromie und starke Charaktersummen sowie die
Reichweite des klassischen $(3,4,5)$-Resultats. Die elementare Ungleichung
$\Xi\ge\log4$ liefert keine IUT-/gcd-Brücke; historische Diagnosescores und
fest gesetzte Gates liefern keinen Beal-Ausschluss. Numerische Tabellen
bleiben historische Diagnostik. HD3', getypte Transfers, explizite Konstanten,
unabhängige Methodenvalidierung und vollständige Quellen-/Beweisprüfung bleiben offen.
\end{abstract}

\noindent\textbf{Schlüsselwörter:} Bealsche Vermutung, Height Dominance, Radikal, Potenzgleichungen, primitive Lösungen, Mikrocluster-Architektur, Potenzrest-Sieb, großes Sieb, Charaktersummen.

\clearpage
\tableofcontents
\clearpage

\bigskip
\hrule
\medskip
\begin{center}
\textbf{Teil A. Height-Dominance Architektur (Deterministische Route)}
\end{center}
\medskip
\hrule
\bigskip

\section{Zweck}

Die Bealsche Vermutung~\cite{Mauldin1997} besagt, dass jede Lösung von
\[
  A^x + B^y = C^z,
  \qquad A,B,C,x,y,z \in \mathbb Z_{>0},
  \qquad x,y,z>2,
\]
einen gemeinsamen Primteiler von $A,B,C$ besitzt. Äquivalent dazu gibt es keine primitive Lösung mit $\gcd(A,B,C)=1$.

Teil A importiert keinen Beweis aus der Riemannschen Vermutung. Wir übernehmen lediglich eine Beweisarchitektur: den relevanten Cluster intrinsisch zu definieren, die Aufhebung (Cancellation) zu isolieren und die verbleibende Masse oder den Defekt in ein zwingendes Komplement zu verschieben. Die modulare Gesamtstruktur dieses Reduktions- und Verifikationsprogramms ist in \cref{fig:bhd_architecture} zusammenfassend dargestellt; der logische Kern der drei Reduktions-Gates ist in \cref{tab:reduction_matrix} tabelliert.

\section{Primitives Potenzgerüst}

\begin{definition}[Primitives Beal-Gegenbeispiel]
Ein primitives Beal-Gegenbeispiel ist ein Tupel $(A,B,C;x,y,z)$ mit $A^x+B^y=C^z$, $x,y,z>2$ und $\gcd(A,B,C)=1$.
\end{definition}

\begin{lemma}[Paarweise Trennung]
\label{lem:pairwise}
In einem primitiven Beal-Gegenbeispiel sind $A,B,C$ paarweise teilerfremd.
\end{lemma}

\begin{proof}
Wenn eine Primzahl $p$ sowohl $A$ als auch $B$ teilt, dann teilt $p$ auch $A^x+B^y=C^z$, folglich $p \mid C$, was der Primitivität widerspricht. Das gleiche Argument gilt für die Paare $(A,C)$ und $(B,C)$.
\end{proof}

\begin{definition}[Höhe und Radikal einer Lösung]
Für ein primitives Tupel setzen wir
\[
  H_* := C^z, \qquad R := \rad(ABC), \qquad
  \sigma := \frac1x+\frac1y+\frac1z.
\]
\end{definition}

\begin{remark}[Notation]
In Teil A verwenden wir $H_*$ für die Höhe einer einzelnen Lösung, um sie optisch von der globalen Höhenschranke $H$ zu unterscheiden, die in der Zählfunktion $T_K(H)$ von Teil B auftritt.
\end{remark}

Die Exponentensumme $\sigma=1/x+1/y+1/z$ liefert die
Orbifold-Euler-Charakteristik $\chi=\sigma-1$. Ihr Vorzeichen klassifiziert
Signaturen, liefert jedoch keine exakte Geschlechtsformel für eine
nicht konstruierte Überlagerung.
Für eine tatsächlich konstruierte glatte Galoisüberlagerung von
$\mathbb P^1$ vom Grad $d$ mit genau den angegebenen drei Trägheitsordnungen
gilt nach Riemann--Hurwitz $2g-2=d(-\chi)$, also im hyperbolischen Fall
$g=1+d|\chi|/2$. Die Setzung $d=\mathrm{lcm}(x,y,z)$ ohne Konstruktion
ist unbegründet: $(3,3,4)$ ergäbe sogar $g=3/2$.
Der externe Darmon--Granville-Endlichkeitssatz für feste nichtverschwindende
Koeffizienten und feste hyperbolische Signaturen bleibt davon getrennt
\cite{DarmonGranville1995}; er liefert weder Leerheit noch einen uniformen
Exponentensatz. In der nichtkubischen Beal-Domäne gilt $\chi\le-1/12$.
\begin{table}[!htbp]
\centering\small
\begin{tabularx}{\textwidth}{@{}l c >{\raggedright\arraybackslash}p{3.1cm} >{\raggedright\arraybackslash}X@{}}
\toprule
Regime & $\chi$ & Beispiele & Zulässige Einordnung\\
\midrule
Sphärisch & $>0$ & $(2,2,z)$, $(2,3,5)$ & Signaturklasse; hier kein uniformer Lösungszähler hergeleitet.\\
\addlinespace
Parabolisch & $=0$ & $(2,3,6)$, $(2,4,4)$, $(3,3,3)$ & Kubischer Fermat separat; Modelle und Koeffizienten sind wesentlich.\\
\addlinespace
Hyperbolisch, gemischt & $<0$ & $(2,3,7)$, $(2,4,5)$ & Externe Endlichkeit bei fester Signatur, keine effektive Leerheit.\\
\addlinespace
Hyperbolisch, Beal & $\le-1/12$ & $x,y,z\ge3$, ohne $(3,3,3)$ & Beal bleibt offen; Krümmungsvorzeichen ist kein Ausschlussbeweis.\\
\bottomrule
\end{tabularx}
\caption{Orbifold-Signaturklassifikation. Insbesondere ist $(2,2,z)$ wegen $\chi=1/z>0$ nicht hyperbolisch. Das Geschlecht einer Überlagerung benötigt ihren tatsächlichen Grad und ihre Verzweigung.}
\label{tab:regime_matrix}
\end{table}

\section{Potenz-Kompression}

\begin{proposition}[Potenz-Kompression, HD-PC]
\label{prop:compression}
Jedes primitive Beal-Gegenbeispiel erfüllt
\[
  R \le ABC < H_*^{\sigma}.
\]
\end{proposition}

\begin{proof}
Da $A^x < C^z=H_*$ und $B^y < C^z=H_*$, haben wir $A < H_*^{1/x}$, $B < H_*^{1/y}$, $C=H_*^{1/z}$. Somit ist $ABC < H_*^{1/x+1/y+1/z}=H_*^{\sigma}$. Da $R=\rad(ABC)\le ABC$, folgt die Behauptung.
\end{proof}

\begin{remark}[Exponenten-Lücke]
Das Maximum von $\sigma$ für $x,y,z>2$ ist $1$, das nur bei $(x,y,z)=(3,3,3)$ erreicht wird. Dieser Fall ist durch den Fermatschen Satz für den Exponenten 3 ausgeschlossen. Für jedes verbleibende Exponententripel hat man $\sigma<1$.
\end{remark}

\section{Potenz-Radikal-Mikrocluster}

\begin{definition}[Potenz-Radikal-Mikrocluster]
Seien $x,y,z>2$ fest vorgegeben. Der Potenz-Radikal-Mikrocluster $V_{HD}(x,y,z)$ ist die Datenstruktur von drei disjunkten Primzahlpaketen $P_A, P_B, P_C$ und Bewertungsvektoren, die auf diesen Paketen getragen werden und $A^x, B^y, C^z$ erzeugen, wobei die Exponenten durch die Vielfachen $x\mathbb Z$, $y\mathbb Z$ und $z\mathbb Z$ entsprechend eingeschränkt sind.
\end{definition}

\begin{remark}[Nicht-Zirkularität]
Der Cluster wird ausschließlich aus der Primitivität und den Exponentendaten definiert. Er setzt nicht voraus, dass sich die additive Gleichung schließen kann oder nicht. Dies ist die Aussage des Endspiels.
\end{remark}

\section{Cancellation-Koordinaten}

Das Teilen von $A^x+B^y=C^z$ durch $H_*=C^z$ liefert
\[
  \frac{A^x}{H_*}+\frac{B^y}{H_*}=1.
\]
Dies ist die skalare Höhen-Aufhebung (scalar height cancellation): Der führende Höhenterm wird exakt durch die beiden eingehenden Potenzterme ausgeglichen.

\begin{definition}[Höhendefekt]
Der Höhendefekt eines primitiven Tupels ist
\[
  D_{HD} := \frac{H_*^{\sigma}}{R}.
\]
Nach \Cref{prop:compression} erfüllt jedes primitive Gegenbeispiel $D_{HD}>1$.
\end{definition}

Der gewünschte Höhendominanz-Mechanismus (Height Dominance) muss die entgegengesetzte Ungleichung erzwingen oder stattdessen einen gemeinsamen Primteiler bedingen.

\section{Das Drei-Lemma-Endspiel}

\begin{lemma}[HD1: Skalare Höhen-Aufhebung]
\label{lem:hd1}
Jede primitive Lösung erfüllt
$A^x/H_* + B^y/H_* = 1$.
\end{lemma}

\begin{proof}
Dies ist die ursprüngliche Gleichung geteilt durch $H_*=C^z$.
\end{proof}

\begin{conjecture}[HD2': Formales Verifier-Entropiebudget (Teil-B-Form)]
\label{conj:hd2prime}
Für jedes feste $K \ge 4$ existieren Konstanten $A=A(K) \ge 1$,
$c=c(K) > 0$ und eine Schwelle $H_{0}=H_{0}(K)$, sodass für jeden
primitiven Potenz-Radikal-Mikrocluster mit Exponenten $3 \le x,y,z \le K$
und $H_* \ge H_{0}$ die polylogarithmische Verifier-Skala
$Q_{\min} = (\log H_*)^{A}$ erfüllt:
\[
  I(Q_{\min};x,y,z) \ \ge\ c \log H_*,
\]
wobei $I$ die kumulative Verifier-Entropie aus \Cref{def:Iverifier} ist.
Insbesondere wächst das Verifier-Budget, das nötig ist, um die Höhe
$H_*$ zu dominieren, polylogarithmisch in $H_*$, uniform in zulässigen
Exponententripeln.
\end{conjecture}

\begin{remark}[Korrigierter Status des formalen HD2'-Budgets]
Das historische Vermutungslabel bleibt als Verweis erhalten; die angegebene
Existenzform folgt jedoch asymptotisch aus lokalen Dichten und Primzahlprogressionen
bei festem $K$: $I(Q;x,y,z)\gg_K Q/\log Q$, uniform über die endliche
Exponentenbox, und $Q=(\log H_*)^2$ genügt für große Höhe.
Dieses Budget ist von $A,B,C$ unabhängig und zertifiziert weder Wurzelsektion
noch Beyond-CRT-Transfer, explizite Höhenschwelle oder HD3'.
Die 1000 Near-Misses sind historische Diagnostik, kein Signifikanztest und
keine unabhängige Validierung der lokal-globalen Methode.
\end{remark}

\begin{conjecture}[HD3': Zwingendes Verifier-Komplement -- Beal-starke
lokal-globale Form]
\label{conj:hd3prime}
Wenn ein primitives Tupel in $V_{HD}(x,y,z)$ $V_q(A,B;x,y,z)=1$ für
jede Primzahl $q$ erfüllt, dann ist $\gcd(A,B,C)>1$.

Jeder exakte positive primitive Abschluss erfüllt alle notwendigen
Verifier-Bedingungen. Die angegebene HD3'-Implikation würde daher, wenn
sie mit einem wohldefinierten Tupel-/Höhenvertrag bewiesen wird, Beal
implizieren. Keine strikte Stärker-/Schwächerbeziehung zur endlichen
Root-Coherence-Variante ist hier bewiesen.
Die Pass/Fail-Bedingung vergisst die Wurzel und bindet das sonst freie $C$
nicht; Wurzelidentität, Höhe und Primzahlquantoren benötigen einen eigenen
arithmetischen Vertrag. Diese Vermutung ist nicht bewiesen.
\end{conjecture}
\begin{remark}[Status von HD3']
HD1 und HD-PC sind Normalisierungen; HD3' trägt den offenen arithmetischen
Ausschluss. Der bedingte Abschluss ist eine formale Reduktion, kein Beal-Beweis.
Ein externer Satz über exakte $(3,4,5)$-Lösungen beweist HD3' nicht für
jedes lokal passierende Tupel mit ungebundenem $C$.
\end{remark}

\begin{table}[htb]
\centering
\small
\begin{tabularx}{\textwidth}{@{} >{\raggedright\arraybackslash}p{0.15\linewidth} >{\raggedright\arraybackslash}p{0.25\linewidth} >{\raggedright\arraybackslash}p{0.25\linewidth} >{\raggedright\arraybackslash}X @{}}
\toprule
Reduktions-Gate & Mathematische Aussage & Beweis- / Statuslevel & Diagnostik \& Guardrails \\
\midrule
HD1 (Skalare Höhen-Aufhebung) & $A^x/H_* + B^y/H_* = 1$ für primitive Lösungen & Bewiesen (algebraische Identität) & Strikte Begrenzung auf primitive Tripel; Null-Dominanz-Referenzpunkt \\
\addlinespace
HD2' (Formales Entropiebudget) & $I((\log H_*)^2;x,y,z)\ge c_K\log H_*$ asymptotisch & Folgt aus lokalen Eingaben bei festem $K$ & Basenunabhängig; keine Wurzelkohärenz zertifiziert \\
\addlinespace
HD3' (Arithmetischer Ausschluss) & Getypte Wurzel-/Höhen-/Primzahlbrücke zum Abschluss & Offene Hypothese & Diagnostik und synthetische Kontrollen liefern diese Brücke nicht \\
\bottomrule
\end{tabularx}
\caption{Die algebraische HD1-Normalisierung, das formale asymptotische HD2'-Entropiebudget und die offene arithmetische HD3'-Brücke haben verschiedene Rollen. Historische Diagnostik validiert die Brücke nicht.}
\label{tab:reduction_matrix}
\end{table}

\begin{figure}[htbp]
\centering
\resizebox{\textwidth}{!}{%
\begin{tikzpicture}[
 box/.style={rectangle,draw,rounded corners=3pt,align=center,
             text width=5.4cm,inner sep=7pt,font=\small},
 proven/.style={box,fill=blue!8,draw=blue!70},
 budget/.style={box,fill=teal!8,draw=teal!70},
 openbarrier/.style={box,fill=orange!8,draw=orange!70},
 closure/.style={box,fill=purple!8,draw=purple!70},
 arrow/.style={-{Stealth[length=6pt]},thick,draw=black!70},
 dasharrow/.style={arrow,dashed}
]
\node[box,text width=11.5cm,fill=gray!10] (input) at (0,1.3) {\textbf{Hypothetisches primitives Gegenbeispiel}\\$A^x+B^y=C^z$, $\gcd(A,B,C)=1$, $x,y,z\ge3$};
\node[align=center,font=\small\bfseries] at (-3.5,-0.7) {TEIL A};
\node[align=center,font=\small\bfseries] at (3.5,-0.7) {TEIL B};
\node[proven,anchor=north] (hd1) at (-3.5,-1.8) {\textbf{HD1 und HD-PC}\\Skalare Aufhebung und\\$\rad(ABC)<H_*^\sigma$\\Elementare Normalisierungen};
\node[budget,below=0.5cm of hd1] (hd2) {\textbf{Formales HD2'-Budget}\\$I((\log H_*)^2;x,y,z)\ge c_K\log H_*$\\Asymptotischer lokaler Input bei festem $K$\\Keine Wurzelsektion zertifiziert};
\node[openbarrier,below=0.5cm of hd2] (hd3) {\textbf{HD3': Offene arithmetische Brücke}\\Wurzelidentität, tatsächliche Summe, Höhe und Primzahlquantoren müssen gebunden sein\\Diagnostik liefert keinen Ausschluss};
\node[closure,below=0.5cm of hd3] (closureA) {\textbf{Bedingter Abschluss}\\Eine bewiesene getypte HD3'-Brücke würde einen gemeinsamen Primteiler erzwingen\\Beal bleibt offen};
\node[proven,anchor=north] (count) at (3.5,-1.8) {\textbf{Positive exakte Lösungszählung}\\$T_K(H)$, positive Basen, festes $K\ge4$\\Exponenten $3,\ldots,K$, ohne $(3,3,3)$};
\node[proven,below=0.5cm of count] (sieve) {\textbf{Additives Sieb für positive Zeilen}\\Die längere Seite wird gesiebt\\Nenner $1+D_Q(A)$\\Ganzzahlige Nullzeile ausgeschlossen};
\node[proven,below=0.5cm of sieve] (bound) {\textbf{Asymptotische Schranke}\\$T_K(H)\ll_{K,\varepsilon}H^{2/3}/\Sigma(Q;K)$\\Großes $H$, genannter $Q$-Bereich\\Keine numerische Schwelle zertifiziert};
\node[openbarrier,below=0.5cm of bound] (expheur) {\textbf{Exponentielle Exaktzählungs-Vermutung}\\$T_K(H)\le C_KH^{\sigma_K}2^{-c_KI_K(Q)}$\\Würde $T_K\equiv0$ erzwingen\\Form für alle $K$ bereits Beal-stark};
\draw[arrow] (input.west) -- ++(-0.4,0) |- (hd1.west);
\draw[arrow] (input.east) -- ++(0.4,0) |- (count.east);
\draw[arrow] (hd1.south) -- (hd2.north);
\draw[dasharrow] (hd2.south) -- (hd3.north);
\draw[dasharrow] (hd3.south) -- (closureA.north);
\draw[arrow] (count.south) -- (sieve.north);
\draw[arrow] (sieve.south) -- (bound.north);
\draw[dasharrow] (bound.south) -- (expheur.north);
\end{tikzpicture}}
\caption{Bedingte Reduktion (Teil A) und asymptotischer Sieb-Begleiter (Teil B). Das formale Entropiebudget zertifiziert keine Wurzelkohärenz. Die arithmetische Brücke bleibt offen; die polynomiale Zählung liefert keinen deterministischen Ausschluss.}
\label{fig:bhd_architecture}
\end{figure}

\section{Bedingter Abschluss}

\begin{theorem}[Height-Dominance-Abschluss, bedingt]
\label{thm:conditional-closure}
Angenommen, \Cref{conj:hd3prime} gilt für alle Exponententripel $x,y,z>2$, wobei der Fall $(3,3,3)$ durch den Fermatschen Satz für den Exponenten~$3$ abgehandelt wird. Dann gilt die Bealsche Vermutung.

Die Verifier-Paket-Vermutung \Cref{conj:hd2prime} ist in dieser letzten Implikation keine eigene logische Voraussetzung. Ihre Rolle ist diagnostisch und konstruktiv: Sie identifiziert die lokale Information, die ein künftiger Beweis von \Cref{conj:hd3prime} über das CRT-Fenster hinaus transportieren müsste.
\end{theorem}

\begin{proof}
Angenommen, es existiert ein primitives Beal-Gegenbeispiel. Nach \Cref{lem:pairwise,prop:compression} liegt es im Potenz-Radikal-Mikrocluster und erfüllt $R < H_*^{\sigma}$. Da die Gleichung exakt ist, ist für jede Primzahl $q$ der Rest $A^x+B^y \bmod q$ eine $z$-te Potenz, nämlich $C^z \bmod q$. Also gilt $V_q(A,B;x,y,z)=1$ für jedes $q$. \Cref{conj:hd3prime} erzwingt dann einen gemeinsamen Primteiler von $A,B,C$, was der Primitivität widerspricht. Daher existiert kein primitives Gegenbeispiel.
\end{proof}

\section{Ansatz zu HD3': Eine ehrliche Skizze}
\label{sec:hd3-approach}

Die deterministische Route steht oder fällt mit HD3'. Umformuliert: \emph{Wenn ein primitives Tupel $(A,B,C;x,y,z)$ im Potenz-Radikal-Mikrocluster $V_q = 1$ für jede Primzahl $q$ erfüllt, dann ist $\gcd(A,B,C) > 1$.}

\subsection*{Was bekannt ist}

HD1 (\Cref{lem:hd1}) ist die durch $H_*=C^z$ geteilte skalare
Gleichung. HD-PC (\Cref{prop:compression}) ist die elementare
Schranke $\rad(ABC)<H_*^\sigma$. Dies sind Normalisierungen,
kein arithmetischer Ausschluss. Die Existenzform von HD2'
(\Cref{conj:hd2prime}) folgt asymptotisch aus lokalen Dichten und
Primzahlprogressionen bei festem $K$; $Q=(\log H_*)^2$ reicht bei
großer Höhe aus. Ihre Entropie hängt von Exponenten und Primzahlen
ab, nicht von den tatsächlichen Basen. Die historischen Near-Misses
sind Diagnostik und beweisen keine Wurzelkohärenz.

\subsection*{Die harte Barriere}

Ein Pass besagt, dass $A^x+B^y$ modulo $q$ irgendeine $z$-te
Wurzel hat; er identifiziert sie nicht mit dem gewählten $C$.
Ein gültiger globaler Ausschluss benötigt einen präzisen Vertrag
zwischen tatsächlicher Summe, Wurzelidentität, Höhe und Primzahlquantoren.
Statistische Unabhängigkeit der Kongruenztests ist für die kurze
ganzzahlige Suchbox nicht gerechtfertigt. Umgekehrt beweist dieses
Paper keinen allgemeinen Abhängigkeitssatz, der einen gemeinsamen
Primteiler erzwingt. Die polynomiale Schranke aus \Cref{thm:sieve}
liefert kein HD3'. Die exponentielle Vermutung zur kumulativen
Lösungszählung \Cref{conj:exp} würde bereits Beal implizieren,
wenn sie für jede feste Exponentenschranke formuliert wird.

\subsection*{Werkzeuge, die Anwendung finden könnten}

Klassische Siebe liefern die additive Schätzung in Teil~B.
Mittelwertbildung, kurze Charaktersummen, kreisteilungstheoretischer
Abstieg und Frey-/S-Unit-Methoden sind mögliche Werkzeuge, die einen
signaturspezifischen Transfer und bei variierenden Körpern oder
Primzahlmengen eine Uniformitätsprüfung benötigen.
Der Endlichkeitssatz für feste Signaturen von Darmon--Granville
\cite{DarmonGranville1995} liefert kein allgemeines Leerheitszertifikat.
Eine Liste möglicher Methoden ist weder ein Anwendbarkeitssatz noch
ein Beweis dafür, dass sämtliche möglichen Transfers scheitern.

\subsection*{Konkrete Teilziele und ihr korrigierter Umfang}
\begin{enumerate}[leftmargin=*]
\item \textbf{HD3'-RC / HD3'-A.} Ein endlicher Root-Coherence- oder
Salt-Closing-Vertrag muss positive Basen, den tatsächlichen Wert
$A^x+B^y$, eine gemeinsame beschränkte ganze Wurzel und die
Primzahl-/Höhenquantoren festlegen. Ein großer formaler
Root-Kapazitätsdefekt ist für sich keine deterministische Obstruktion.
Keine strikte Implikationshierarchie und kein uniformer Ausschluss sind zertifiziert.
\item \textbf{HD3'-B: schwacher Input und verworfene uniforme starke Form.}
Für nichttriviales $\eta$, alle multiplikativen Charaktere bei null
durch null fortgesetzt, setze $H_x=\{\chi:\chi^x=1\}$ und
$J(\chi,\psi)=\sum_t\chi(t)\psi(1-t)$. Die exakte Konsolidierung ist
\[
 S_\eta(x,y;q)=(q-1)
 \sum_{\substack{\chi\in H_x\setminus\{1\},\,\psi\in H_y\setminus\{1\}\\
                 \chi\psi=\eta^{-1}}}J(\chi,\psi).
\]
Terme mit einem trivialen Charakter heben die Achsenbeiträge auf;
der Faktor $q-1$ darf nicht entfallen. Klassische Jacobi-Schranken
\cite[Ch.~8]{IrelandRosen} liefern den schwachen $O_K(q^{3/2})$-Input.
Für $q\equiv1\pmod{12}$, $(x,y)=(3,4)$ und $\eta$ der Ordnung $12$
existiert ein Resonanzpaar, also $|S_\eta|=(q-1)\sqrt q$.
Dies widerlegt eine pauschale $\eta$-uniforme $q^{1+\varepsilon}$-Schranke
bei $\varepsilon<1/2$ auch im teilerfremden gemischten Fall.
Eingeschränkte $(3,4,5)$-Aussagen werden dadurch nicht widerlegt.

Die skalare Familie $a^x+b^y=c$ wird nach endlichem étalem Basiswechsel
$c=t^L$ auf $\mathbb G_m$, $L=\mathrm{lcm}(x,y)$, konstant:
$a=t^{L/x}u$, $b=t^{L/y}v$ reduziert auf $u^x+v^y=1$.
In Charakteristik ohne Teiler von $xy$ faktorisiert ihre geometrische
Monodromie durch eine endliche Deckgruppe; diese konkrete Familie
kann keine positivdimensionale volle symplektische Monodromiegruppe besitzen.
Keine Trivialisierung bei $c=0$, kein allgemeines Frey-Familien-No-Go
und keine Xu-Äquivalenz \cite{Xu2020,FouvryKatz2001,KatzTwistedL} folgen.
\item \textbf{HD3'-C / HD3'-D.} Zählungen aller lokal passierenden Tupel
und exakter positiver Lösungen sind verschieden. Keine Äquivalenz
dieser Zähl-/Salt-/Kofaktorziele mit \cref{conj:exp} ist bewiesen.
\item \textbf{HD3'-E: externer Satz über exakte Lösungen.}
Siksek--Stolls Theorem 1.1 \cite{SiksekStoll2012} bestimmt für
$u^3+v^4+w^5=0$ nur triviale primitive ganzzahlige Lösungen.
Mit $w=-C$ schließt dies positive primitive $A^3+B^4=C^5$ aus.
Es beweist weder die vorgeschlagene Verifier-Aussage für jedes lokal
passierende Tupel mit ungebundenem $C$ noch validiert es unsere Architektur.
Das korrigierte asymptotische Sieb bei $K=5$, $Q=H^{1/10}$ liefert nur
$T_{(3,4,5)}(H)\ll H^{29/60}\log H$; $Q=H^{1/5}$ liegt außerhalb
der sauberen uniformen Reichweite.
\end{enumerate}

\subsection*{Was offen bleibt}
HD3' und ein getypter Vergleich niedriger Wurzeln bleiben offen.
Das klassische additive Sieb liefert die vorgeschlagene exponentielle
kumulative Zählschranke nicht. Letztere ist unter den formalen
Entropie-Eingaben bereits Beal-stark, siehe \cref{conj:exp}.
Fehlgeschlagene Kandidaten oder eingeschränkte Quellensätze beweisen
keinen Unmöglichkeitssatz für sämtliche Methoden.

\section{Was nicht verwendet wird}

Keine abc-, Szpiro- oder Riemann-Implikation wird zum Abschluss von
Beal verwendet. Der Zookeeper-Vergleich ist architektonisch.
Eine gemeinsame geometrische oder Operator-Umgebung für einen
Beal-Transfer wurde hier nicht konstruiert; ihr Fehlen ist kein
allgemeiner Unmöglichkeitssatz für einen Hilbert--Pólya-Weg oder
für gemeinsame Geometrie.

\bigskip\hrule\medskip
\begin{center}
\textbf{Teil B. Potenzrest-Sieb-Begleiter (Asymptotische Schranke)}
\end{center}
\medskip\hrule\bigskip

\section{Das Sieb-Setup}
\label{sec:weilinput}

Teil~B liefert auf benannten klassischen Inputs eine uniforme
asymptotische Zählschätzung für die endliche Exponentenbox.
Sie ergänzt den qualitativen Endlichkeitssatz für feste Signaturen
\cite{DarmonGranville1995}, zertifiziert aber keine explizite
Höhenschwelle. Die Entropietabelle ist historische beschreibende
Evidenz. Die exponentielle Vermutung zur positiven Lösungszählung
ist bereits Beal-stark; sie ist keine etablierte schwächere Brücke
zu HD3'.

\section{Potenzrest-Verifier (Teil B)}

Sei $q$ eine beliebige Primzahl und $z \ge 3$. Setze
\[
  R_z(q) = \{u^z \bmod q : u \in \mathbb{F}_q\},
  \qquad
  |R_z(q)| = 1 + (q-1)/\gcd(z,q-1).
\]

\begin{definition}[Lokale Akzeptanzdichte]
\label{def:delta}
Für eine Primzahl $q$ mit $\gcd(z,q-1)>1$ und Exponenten $x,y,z>2$,
\[
  \delta_q(x,y,z) :=
    \frac{|\{(a,b)\in \mathbb{F}_q^2 : a^x+b^y \in R_z(q)\}|}{q^2}.
\]
Für triviale $q$ (d.h.\ $\gcd(z,q-1)=1$) liefert das Sieb keine Information, Dichte- und Entropiesummen verwenden daher nur nichttriviale Primzahlen.
\end{definition}

\begin{definition}[Verifier-Indikator]
\label{def:Vq}
Für ein Tupel $(A,B,x,y,z)$ und jede Primzahl $q$,
\[
  V_q(A,B;x,y,z) :=
    \mathbf{1}\bigl\{(A^x+B^y) \bmod q \in R_z(q)\bigr\}.
\]
Bei trivialen Primzahlen $\gcd(z,q-1)=1$ gilt
$R_z(q)=\mathbb F_q$ und $V_q=1$. Dichtesummen verwenden nur
nichttriviale Primzahlen.
\end{definition}

\begin{lemma}[Notwendige Bedingung]
\label{lem:necessary}
Wenn $(A,B,C;x,y,z)$ $A^x+B^y=C^z$ erfüllt, dann gilt
\[
  V_q(A,B;x,y,z)=1
\]
für jede Primzahl $q$.
\end{lemma}

\begin{proof}
Es gilt $A^x+B^y \equiv C^z \pmod q$, und $C^z \bmod q \in R_z(q)$ nach Definition.
\end{proof}

\begin{definition}[Kumulative Verifier-Entropie]
\label{def:Iverifier}
Für eine Verifier-Schranke $Q$ und Exponenten $x,y,z>2$,
\[
  I(Q;x,y,z) := \sum_{q \le Q,\ \gcd(z,q-1)>1} -\log_2 \delta_q(x,y,z).
\]
\end{definition}

\section{Worst-Case-Akzeptanzdichte}

Sei $K\ge4$ eine feste ganze Zahl. Definiere für jede Primzahl mit
nichtleerer zulässiger Indexmenge
\[
 \delta_{q,K}^{*}=\max_{\substack{3\le x,y,z\le K\\
                         \gcd(z,q-1)>1}}\delta_q(x,y,z).
\]
Primzahlen mit leerer Indexmenge werden aus Summen mit dieser Größe
weggelassen; einem leeren Maximum wird kein Zahlenwert zugewiesen.

\begin{lemma}[Uniforme Charaktersummen-Schranke für $\delta_q$]
\label{lem:weil}
Für $3\le x,y,z\le K$, $q>K^4$, $q\nmid xyz$ und
$d=\gcd(z,q-1)>1$ liefert der klassische Jacobi-Summen-Input
\[
 \left|\delta_q(x,y,z)-\frac1d\right|\le\frac{2K}{\sqrt q}.
\]
\end{lemma}

\begin{proof}
Schreibe $H_m=\{\chi:\chi^m=1\}$ für multiplikative Charaktere von
$\mathbb F_q^*$. Jeder Charakter wird bei null durch null fortgesetzt,
auch der triviale Charakter. Setze $N_0=\#\{(a,b):a^x+b^y=0\}$
und $S_\eta=\sum_{a,b}\eta(a^x+b^y)$. Der Potenzrest-Indikator ergibt
die exakte Identität
\[
 q^2\delta_q=\frac{q^2}{d}
       +N_0(1-1/d)+\frac1d\sum_{\substack{\eta\in H_z\\\eta\ne1}}S_\eta.
\]
Bei $a=0$ trägt nur $b=0$ zu $N_0$ bei; bei $a\ne0$ gibt es
höchstens $K$ Möglichkeiten für $b$. Also $N_0\le1+K(q-1)\le Kq$.

Für nichttriviales $\eta$ liefern Faltung und Charakterorthogonalität
\[
 S_\eta=(q-1)
 \sum_{\substack{\chi\in H_x\setminus\{1\},\
                 \psi\in H_y\setminus\{1\}\\
                 \chi\psi=\eta^{-1}}}J(\chi,\psi).
\]
Die Faserzählung der Potenzabbildung ist nämlich
$\sum_{\chi\in H_x}\chi(u)+\mathbf1_{\{u=0\}}$.
Für $u+v=s\ne0$ lässt die Substitution $u=st$, $v=s(1-t)$
die Summe $\sum_{s\ne0}(\chi\psi\eta)(s)$ übrig, die außerhalb
der Resonanz verschwindet. Resonante Terme mit einem trivialen
Charakter haben Jacobi-Summe $-1$; sie heben sich exakt gegen die
entsprechenden Achsenbeiträge auf. Das Paar mit zwei trivialen
Charakteren ist wegen $\eta\ne1$ nicht resonant.

Es verbleiben höchstens $K$ resonante Paare, deren Produkt
$\chi\psi=\eta^{-1}$ nichttrivial ist. Die klassische Jacobi-Summen-
Identität liefert $|J(\chi,\psi)|=\sqrt q$
\cite[Kap.~8]{IrelandRosen}. Daher gilt
$|S_\eta|\le K(q-1)\sqrt q$. Nach Division der exakten Identität
durch $q^2$ ist der Fehler höchstens
$K/q+K/\sqrt q\le2K/\sqrt q$.
\end{proof}

\begin{remark}[Umfang des Inputs]
Das Argument verwendet die benannte klassische Jacobi-Summen-
Identität mit den oben angegebenen Charakter- und Randkonventionen.
Es liefert eine lokale Dichteschätzung. Es liefert weder eine uniforme
$q^{1+\varepsilon}$-Schranke für alle verdrehten zweidimensionalen
Summen noch einen globalen Ausschlusssatz. Historische Begleitnotizen
und ihre Prüfsiegel bleiben als historische Nachweise erhalten.
\end{remark}

\begin{lemma}[Worst-Case obere Schranke]
\label{lem:worstcase}
Für zulässige Primzahlen $q>K^4$ gilt
\[
 \delta_{q,K}^{*}\le\frac12+\frac{2K}{\sqrt q}.
\]
\end{lemma}
\begin{proof}
Für jedes zulässige Tripel gilt $d\ge2$, also $1/d\le1/2$.
Wende \Cref{lem:weil} an und bilde anschließend das Maximum.
\end{proof}

\section{Die asymptotische Siebschranke}
\label{sec:sieve}

\begin{theorem}[Uniforme asymptotische Siebschranke]
\label{thm:sieve}
Seien $K\ge4$ eine feste ganze Zahl und $\varepsilon>0$.
Die Funktion $T_K(H)$ zählt primitive Tupel mit
\emph{positiven ganzzahligen} Basen $A,B,C$, $A^x+B^y=C^z$,
$3\le x,y,z\le K$, $C^z\le H$ und $(x,y,z)\ne(3,3,3)$. Definiere
\[
 \Sigma(Q;K)=
 \sum_{\substack{q\le Q\\
      \exists\,3\le z'\le K:\ \gcd(z',q-1)>1}}
       \frac{1-\delta_{q,K}^{*}}{\delta_{q,K}^{*}}.
\]
Unter den unten angegebenen klassischen Inputs zu Charaktersummen,
additivem Großsieb und dem Primzahlsatz für feste Restklassen gilt
für hinreichend großes $H$ und jedes
\[
 (\log H)^{2+\varepsilon}\le Q\le H^{1/(2K)}
\]
die Schranke
\[
 T_K(H)\ll_{K,\varepsilon}\frac{H^{2/3}}{\Sigma(Q;K)}.
\]
Die untere Grenze ist hinreichend, nicht notwendig. Diese Aussage
zertifiziert weder explizite implizite Konstanten noch eine numerische
Schwelle $H_0(K,\varepsilon)$ oder die behauptete Ordnung
$K^{\Theta(K)}$.
\end{theorem}

\begin{remark}[Definitionsbereich und Parameterbereich]
Für feste $K,\varepsilon$ ist der angegebene Bereich bei großem $H$
nichtleer. Für kleine Höhen gilt die triviale Schranke aus der
positiven Suchbox; sie werden nicht stillschweigend in den
asymptotischen Vergleich einbezogen. Für $K=3$ wird das einzige Tripel
durch den klassischen kubischen Fermat-Fall entfernt, und $T_3=0$;
es wird kein leeres Maximum verwendet. Die uniforme obere Grenze
absorbiert $Q^2$ in die längere Seite der Suchbox.
\end{remark}

\begin{lemma}[Additives Großsieb für positive Zeilen]
\label{lem:LS2}
Bei $M$ positiven Zeilen laufe die Variable jeder Zeile über ein
Intervall aus $L$ aufeinanderfolgenden ganzen Zahlen. Für jede
Primzahl $q\le Q$ haben die zulässigen Zeilenklassen Dichte
$\delta_q(A)$; setze
\[
 D_Q(A)=\sum_{q\le Q}\frac{1-\delta_q(A)}{\delta_q(A)}.
\]
Ist eine zulässige Menge leer, hat die betreffende Zeile keine
Survivors. Andernfalls ist die Gesamtzahl der Survivors höchstens
\[
 \sum_{A=1}^{M}\min\left\{L,\
       \frac{L-1+Q^2}{1+D_Q(A)}\right\}.
\]
Der Ausdruck ist insbesondere auch bei $D_Q(A)=0$ definiert.
\end{lemma}

\begin{proof}
Fixiere eine Zeile und bezeichne ihre Survivor-Anzahl mit $Z$.
Schreibe $S(\alpha)=\sum_{n\ {\rm surviving}}e^{2\pi i n\alpha}$.
Parseval und Cauchy--Schwarz auf den zulässigen Restklassen liefern
\[
 \sum_{a=1}^{q-1}|S(a/q)|^2
       \ge Z^2\frac{1-\delta_q(A)}{\delta_q(A)}.
\]
Die Frequenz null und die gekürzten Brüche $a/q$ für verschiedene
Primzahlen sind modulo eins verschieden und $Q^{-2}$-separiert.
Das klassische additive Großsieb
\cite{Bombieri1971,KedlayaANT} ergibt daher
\[
 Z^2(1+D_Q(A))\le(L-1+Q^2)Z.
\]
Für $Z>0$ dividiere durch $Z$; für $Z=0$ folgt die Aussage sofort.
Außerdem gilt $Z\le L$. Summiere über die positiven Zeilen.
Der genaue additive Input steht in
\cite[Satz~18.1]{KedlayaANT}.
\end{proof}

\begin{remark}[Zeilendichten und Vergleich mit $\Sigma$]
\label{rem:Sigma-comparison}
Setze für die Klassen $a^x+b^y\in R_z(q)$
$d_y=\gcd(y,q-1)$ und $d_z=\gcd(z,q-1)$.
Die Zeile mit Rest null hat die exakte Dichte
\[
 \delta_q(0)=
 \frac{1+(q-1)\gcd(d_y,d_z)/d_z}{q}.
\]
Somit gilt $\delta_q(0)=1$ genau dann, wenn $d_z\mid d_y$.
Beispielsweise ergeben $q=7,y=3,z=6$ den Wert $4/7$,
dagegen $y=6,z=3$ den Wert $1$.

Für $a\ne0$ und $q\nmid y$ ist das Polynom $a^x+b^y$ quadratfrei.
Die klassische eindimensionale multiplikative Charakterschranke gibt
\[
 \left|\delta_q(a)-\frac1{d_z}\right|\le2Kq^{-1/2}.
\]
Die Nullstellen des Polynoms tragen dabei höchstens $K/q$ bei,
die nichttrivialen Charaktersummen höchstens $K/\sqrt q$.
Folglich gelten bei \emph{nichttrivialen} Primzahlen $d_z>1$ und
$q>144K^2$ die Schranke $\delta_q(a)\le2/3$ und das Siebgewicht
mindestens $1/2$. Bei trivialen Primzahlen $d_z=1$ sind die Dichte
$1$ und das Gewicht null; diese Primzahlen dürfen nicht für die
untere Schranke verwendet werden.

Für eine positive ganzzahlige Zeile $A$ gilt $A\bmod q=0$ nur
bei ihren Primteilern. Mit
$\pi_z(Q)=\#\{q\le Q:\gcd(z,q-1)>1\}$ folgt
\[
 D_Q(A)\ge
 \tfrac12\{\pi_z(Q)-\omega(A)-\pi(144K^2)\}.
\]
Die Zeile $A=0$ ist ausgeschlossen; $\omega(0)$ ist nicht definiert.
Die elementare Schranke $\omega(A)\le\log_2 H$ für $1\le A\le H$
reicht hier aus, ist aber asymptotisch nicht scharf:
Aus $\omega(A)=k$ folgt $A\ge(k+1)!$ und damit
$\omega(A)=O(\log H/\log\log H)$.

Für festes $3\le z\le K$ sind Primzahlen $q\equiv1\pmod{\ell}$
für einen Primteiler $\ell$ von $z$ nichttrivial für $z$.
Der Primzahlsatz für feste arithmetische Progressionen
\cite[Satz~4.12]{KedlayaANT} liefert
$\pi_z(Q)\asymp_K Q/\log Q$ für $Q\to\infty$.
Der Zeilennenner verwendet diese Primzahlanzahl für ein
\emph{festes $z$}, nicht die größere Vereinigung zulässiger
Primzahlen für verschiedene $z'$.
Bei $K\ge4$ macht $z'=4$ jede ungerade Primzahl zulässig, und
\Cref{lem:weil,lem:worstcase} liefern
$\Sigma(Q;K)\asymp_K Q/\log Q$: Bei großen Primzahlen liegt jede
lokale Dichte zwischen einer positiven, von $K$ abhängigen Konstante
und $2/3$; die endlich vielen kleineren Primzahlen haben endliche
Gewichte.

Die hinreichende untere Grenze aus \Cref{thm:sieve} lässt daher
$\pi_z(Q)$ bei hinreichend großem $H$ die Größe
$\log_2 H+\pi(144K^2)$ uniform in den positiven Zeilen dominieren.
Die obere Grenze $H^{1/(2K)}$ wird nicht anstelle des tatsächlichen
$Q$ eingesetzt. Die untere Grenze ist weder notwendig noch eine
Äquivalenz. Wächst $Q$ unterhalb dieser Grenze, muss $\Sigma(Q;K)$
nicht beschränkt bleiben; eine gemittelte Erweiterung auf diesen
Bereich ist hier nicht nachgewiesen.
\end{remark}

\begin{proof}[Beweis von \Cref{thm:sieve}]
Fixiere ein zulässiges Exponententripel und setze
$U=H^{1/\min(x,y)}$, $V=H^{1/\max(x,y)}$.
Die positive $(A,B)$-Box hat eine längere Seite von höchstens $U$
und eine kürzere von höchstens $V$. Vertausche die Variablen bei
Bedarf: Die kürzere Seite indiziert die positiven Zeilen, die längere
wird gesiebt. Das Zeilendichte-Argument aus
\Cref{rem:Sigma-comparison} gilt mit passend vertauschten $x,y$.

Im gesamten angegebenen Bereich hat bei großem $H$ jede positive
Zeile $D_Q(A)\gg_K Q/\log Q\asymp_K\Sigma(Q;K)$.
Außerdem gilt $Q^2\le H^{1/K}\le U$.
Mit $L=\lfloor U\rfloor$ und $M=\lfloor V\rfloor$ in
\Cref{lem:LS2} ist die Anzahl der Survivor-Paare
\[
 \ll_K\frac{UV}{\Sigma(Q;K)}
       =\frac{H^{1/x+1/y}}{\Sigma(Q;K)}.
\]
Jedes positive Paar hat höchstens eine positive ganze Zahl $C$ mit
$C^z=A^x+B^y$. Die Einschränkung auf primitive Paare und auf
$C^z\le H$ kann die Anzahl nur verringern. Wegen $1/x+1/y\le2/3$
liefert die Summe über endlich viele Exponententripel die angegebene
$\ll_{K,\varepsilon}$-Schranke. Alle impliziten Konstanten bleiben
erhalten; der Beweis ersetzt eine asymptotische Ungleichung nicht
durch eine exakte Ungleichung mit Koeffizient eins.
\end{proof}

\begin{remark}[Ausnahmeprimzahlen und nicht absorbierter Rest]
\label{rem:exceptional}
Für festes $K$ entfernt das Weglassen der endlich vielen kleinen
Primzahlen oder Primteiler von $xyz$ insgesamt $O_K(1)$ Gewicht;
eine Behauptung, jedes einzelne Gewicht sei höchstens eins,
wird nicht benötigt. Für ein festes Tripel ergibt sich unter
derselben hinreichenden unteren Grenze und Großhöhenbedingung
die nicht absorbierte Zeilenschranke
\[
 T_{(x,y,z)}(H)\ll_K
 \frac{H^{1/x+1/y}+H^{1/\max(x,y)}Q^2}{\Sigma(Q;K)}.
\]
Die uniforme obere Grenze im Satz absorbiert den zweiten Term.
Explizite Konstanten, eine zertifizierte numerische Höhenschwelle und
die Erweiterung auf kleinere wachsende Siebbereiche bleiben getrennte
Aufgaben.
\end{remark}

\section{Empirische Verifikation}
\label{sec:empirical}

Die historische Tabelle zu 1000 Near-Miss-Zeilen berichtet vier
verschachtelte Verifier-Schranken. Ihre Werte werden als beschreibende
Diagnostik übernommen; die lokale Dichteasymptotik ist ein getrennter
mathematischer Input.

\begin{table}[htb]
\centering
\begin{tabular}{rrrr}
\toprule
$Q$ & $I(Q)$ Median (bits) & $I(Q)/\pi(Q)$ & Steigung$_{\min}$ vs.\ $\log H_*$ \\
\midrule
$251$  & $38.24$  & $0.71$ & $+0.704$ \\
$503$  & $69.78$  & $0.73$ & $+1.116$ \\
$1009$ & $126.70$ & $0.75$ & $+1.737$ \\
$2003$ & $232.75$ & $0.77$ & $+4.289$ \\
\bottomrule
\end{tabular}
\caption{Historisch berichtete Verifier-Entropie für 1000 Near-Misses mit $\log H_*\in[16,39]$. Die vier verschachtelten Schranken liefern beschreibende Quotienten und Steigungen. Sie beweisen weder Konvergenz oder statistische Signifikanz noch unabhängige Methodenvalidierung.}
\label{tab:emp}
\end{table}

\begin{remark}[Beschreibender Status]
Die Tabelle übernimmt historisch berichtete Werte ohne neuen Lauf.
Das formale Entropiebudget bei festem $K$ folgt aus klassischen
lokalen Inputs und ist von den Basen unabhängig.
Weder die Tabelle noch ihre verschachtelten Primzahlschranken
liefern einen Signifikanztest, eine Wurzelsektion, eine exponentielle
Survivor-Schranke oder HD3'.
\end{remark}

\begin{remark}[Root-Coherence-Guardrail]
Der Verifier $V_q(A,B;x,y,z)$ ist bewusst eine grobe Pass/Fail-Bedingung:
Er fragt, ob modulo $q$ irgendeine lokale $z$-te Wurzel existiert. Eine
exakte Beal-Closure verlangt mehr. Derselbe kleine Integer $C$ muss für
alle betrachteten Primzahlen auf kompatible lokale Wurzeln reduzieren.
Äquivalent: Für
\[
  R_q(A,B)=\{c\in\mathbb F_q:\ c^z\equiv A^x+B^y\pmod q\}
\]
braucht man eine globale Sektion kleiner Höhe $C\le H^{1/z}$ durch die
lokalen Mengen $R_q(A,B)$, nicht nur deren jeweilige Nichtleere. Die
empirische Entropietabelle stützt daher nur die Pass/Fail-Verifier-Schicht.
Eine zukünftige HD3'-Route muss die Wurzelidentität mitführen, etwa über
das in den Begleitnotizen dokumentierte Root-Coherence-Capacity-Framework.
Die S-Unit-/Frey-Shadow- und Q97-plus-Synthetic-Control-Audits vom Juni
2026 verschärfen dieselbe Anforderung: Q97-plus lokales Überleben ist
für sich nie Evidenz. Ein künftiger Survivor darf nur dann stärker
gelesen werden, wenn er Q97-plus ist, aus dem tatsächlichen Wert
$A^x+B^y$ branch-bound bleibt und zusätzlich eine globale Wurzelsektion
kleiner Höhe oder source-bound S-Unit-/Frey-Daten besitzt.
\end{remark}

\begin{table}[htb]
\centering
\small
\begin{tabularx}{\textwidth}{@{} >{\raggedright\arraybackslash}p{0.20\linewidth} >{\raggedright\arraybackslash}p{0.25\linewidth} >{\raggedright\arraybackslash}p{0.22\linewidth} >{\raggedright\arraybackslash}X @{}}
\toprule
Diagnoseebene & Stichprobe und Bereich & Beobachtetes Signal & Einordnung \\
\midrule
Verifier-Entropie & 1000 Near-Misses, $Q\in\{251,503,1009,2003\}$ &
$I(Q)/\pi(Q)=0.71$--$0.77$ &
Historische beschreibende Skalierung; keine Wurzelkohärenz validiert. \\
\addlinespace
Raues Root-Pilotledger & 250 Near-Misses, $Q\in\{251,503,1009\}$ &
$250/250$ leere echte Root-Sets; Exact-Power-Kontrolle $250/250$ &
Lokale Obstruktion tritt vor Root-Drift auf; die Kontrolle prüft die Detektor-Funktion. \\
\addlinespace
Kleiner Survivor-Scan & 250 Near-Misses, $Q\in\{7,13,31,61\}$ &
ein Drift-ohne-Sektion-Fall bei $Q=31$; alle echten Zeilen scheitern lokal bis $Q=61$ &
Die Drift-Zone ist in dieser Stichprobe endlich und übergangsartig. \\
\addlinespace
Above-threshold-Suche & 5000 Near-Misses, Checkpoints $31$--$127$ &
fünf $Q=61$-Pässe ohne niedrige Sektion; $0/5000$ lokale Pässe bei $Q=97$ &
Kein stabiler Above-threshold-Survivor; Folgegates brauchen source-bound Familien und Kontrollen. \\
\addlinespace
S-Unit-/Frey-Shadow-Audit & fünf $Q=61$-Survivors, Juni 2026 &
$0$ Q97-plus-Survivors; $0$ direkte S-Unit-/Frey-Anker &
Es bleibt nur methodischer Shadow-Kontext; keine source-bound Stütze für HD3'. \\
\addlinespace
Q97-plus-Synthetic-Control & natürliche Survivors plus registrierte, verschobene und permutierte Kontrollen &
natürlich $0/5$ Q97-plus; Kontrollen erreichen $Q=127$ nur durch synthetische Sektionen oder Branch-Maps &
Lokale Passaggregate allein sind Negativkontrollen; das Triple-Gate ist verpflichtend. \\
\bottomrule
\end{tabularx}
\caption{Root-Coherence- und Verifier-Diagnostik-Ledger. Die Zeilen
trennen Evidenz für die additive Verifier-Schicht von Tests auf niedrige
Root-Sektionen. Alle Root-Coherence-Zeilen sind interne Diagnostik, keine
Beweisevidenz für HD3'.}
\label{tab:root-diagnostics}
\end{table}

\section{Die stärkere Vermutung}

\begin{definition}[Uniformer Boxexponent und Entropie]
\label{def:sigmaK}
Definiere für eine ganze Zahl $K\ge4$
\[
 \sigma_K=\max_{\substack{3\le x,y,z\le K\\(x,y,z)\ne(3,3,3)}}
                       (1/x+1/y)=2/3,\qquad
 I_K(Q)=\min_{\substack{3\le x,y,z\le K\\(x,y,z)\ne(3,3,3)}}
                       I(Q;x,y,z).
\]
Die Indexmengen sind endlich und nichtleer; kein Maximum für $K=3$
wird verwendet. Der Exponent ist die Obergrenze der konkreten
positiven Boxschätzung aus \Cref{thm:sieve}, kein Obergrenzensatz
für sämtliche arithmetischen Methoden.
Die lokalen Dichten und Inputs zu festen Progressionen implizieren
$I_K(Q)\gg_K Q/\log Q$ für großes $Q$.
\end{definition}

\begin{conjecture}[Exponentielle Schranke der kumulativen Lösungszählung]
\label{conj:exp}
Für jede feste ganze Zahl $K\ge4$ existieren Konstanten $a_K\ge2$,
$c_K>0$ und $C_K>0$, sodass für hinreichend großes $H$ und
$Q=(\log H)^{a_K}$ gilt
\[
 T_K(H)\le C_K H^{\sigma_K}2^{-c_K I_K(Q)}.
\]
Gezählt werden die kumulativen positiven primitiven exakten Lösungen
aus \Cref{thm:sieve}, nicht lokal passierende Paare.
\end{conjecture}

\begin{remark}[Warum dies bereits Beal-stark ist]
Der Logarithmus der rechten Seite ist höchstens
\[
 \sigma_K\log H-c'_K
       \frac{(\log H)^{a_K}}{\log\log H}+O_K(1)\longrightarrow-\infty.
\]
Die Schranke erzwingt daher $T_K(H)<1$ für alle hinreichend großen
$H$. Weil $T_K$ nichtnegativ, ganzzahlig und monoton wachsend ist,
folgt $T_K(H)=0$ bei \emph{jeder} Höhe.
Für alle festen $K$ schließt dies jede positive primitive gemischte
Signatur aus; die kubische Signatur wird klassisch ausgeschlossen.
Umgekehrt setzt Beal alle diese Zählwerte auf null und macht die
angegebene Schranke unmittelbar gültig. Unter den angegebenen
Entropie-Inputs ist die Formulierung für alle $K$ somit Beal-stark.

Ein Dichteprodukt ist auf einem vollständigen CRT-Restklassenbereich
exakt; es ist nicht automatisch eine Survivor-Wahrscheinlichkeit
in einer kurzen Box. Unabhängigkeit jenseits dieses Fensters ist
hier nicht gerechtfertigt. Keine Äquivalenz zu einer Vermutung
über lokal passierende Paare und keine strikte Hierarchie der
Root-Coherence-Teilziele werden behauptet.
\end{remark}

\section{Das CRT-Produktfenster}
\label{sec:crt-window}

Die multiplikative Produktheuristik ist nicht in jedem Bereich falsch. Sie ist im endlichen Fenster des chinesischen Restsatzes exakt, und genau diese Aussage erklärt, wo die Barriere beginnt.

\begin{lemma}[CRT-Produktfenster]
\label{lem:crt-window}
Sei $P$ eine endliche Menge von Primzahlen und $M(P)=\prod_{q\in P}q$.
Für jedes $q\in P$ sei $S_q\subseteq\mathbb F_q^2$ und
$\delta_q=|S_q|/q^2$. Sei $S_M\subseteq(\mathbb Z/M(P)\mathbb Z)^2$
das CRT-Produkt der $S_q$. Dann gilt
\[
  |S_M|=M(P)^2\prod_{q\in P}\delta_q.
\]
Wenn $N_P(X,Y)$ die Paare $(A,B)\in[1,X]\times[1,Y]\cap\mathbb Z^2$
zählt, deren Restklasse modulo $M(P)$ in $S_M$ liegt, dann gilt
\[
  N_P(X,Y)
  = XY\prod_{q\in P}\delta_q
    + O\!\left(\prod_{q\in P}\delta_q\,
      (M(P)X+M(P)Y+M(P)^2)\right).
\]
Insbesondere gelten relative Produktasymptotiken, sobald
$M(P)=o(\min(X,Y))$.
\end{lemma}

\begin{proof}
Der chinesische Restsatz identifiziert $(\mathbb Z/M(P)\mathbb Z)^2$
mit $\prod_{q\in P}\mathbb F_q^2$, was die exakte Kardinalität von
$S_M$ liefert. Jede zulässige Restklasse modulo $M(P)$ besitzt
$XY/M(P)^2+O(X/M(P)+Y/M(P)+1)$ Repräsentanten in der Box. Die Multiplikation
mit $|S_M|$ ergibt die behauptete Formel.
\end{proof}

Für die Beal-Verifier-Boxen hat man $X=H^{1/x}$ und $Y=H^{1/y}$.
Das direkte CRT-Produktfenster verlangt also
\[
  M(P)\ll H^{\min(1/x,1/y)}.
\]
Für $P=\{q\le Q\}$ gilt $\log M(P)=\vartheta(Q)\sim Q$, also reicht
dieses direkte Fenster nur bis
\[
  Q\lesssim \min(1/x,1/y)\log H.
\]
In diesem maximalen Bereich liefert das Produkt der lokalen Dichten
höchstens einen subpolynomialen Gewinn der Form
$\exp(-c\pi(Q))=H^{-c/\log\log H+o(1)}$. Das gewünschte Regime
$Q=(\log H)^A$ mit $A>1$ hat bereits
$M(P)=\exp((\log H)^A)$ und ist damit weit größer als die gewöhnliche
Suchbox.

Folglich ist das Hindernis hinter \Cref{conj:exp} nicht bloß eine
fehlende lokale Charaktersummen"=Verbesserung. Es ist das Problem,
Produktinformation über die CRT"=Fundamentaldomäne hinaus zu
transportieren, etwa durch bilineare Dispersion, Frey"=Familien"=Geometrie,
Selmer- oder Mordell--Weil-Struktur oder ein wirklich neues
Korrelationstheorem.

\section{Vergleich mit Darmon--Granville}

Darmon--Granville \cite{DarmonGranville1995} beweist qualitative
Endlichkeit primitiver Lösungen für jede feste hyperbolische
Signatur. Endlichkeit beweist keine Leerheit.
Teil~B liefert stattdessen die positive Boxzählung
\Cref{thm:sieve}, uniform über die endliche Exponentenbox bei festem
$K$ und hinreichend großer Höhe, auf den angegebenen klassischen
Inputs. Sie liefert keine explizite Konstante oder zertifizierte
numerische Höhenschwelle. Ohne diese zusätzliche Arbeit wird kein
Effektivitätsvorteil gegenüber dem Satz für feste Signaturen behauptet.
Die Signaturklassifikation aus \Cref{tab:regime_matrix} konstruiert
allein keine glatte Überlagerung und bestimmt nicht deren Geschlecht.

\section{Ausblick: Begrenzte Negativtests und offene Transfers}
\label{sec:outlook}

Die folgenden Punkte sind Diagnostik möglicher Methoden, keine
allgemeinen Unmöglichkeitssätze und kein unabhängiger Beal-Beweis.

\subsection{Starke verdrehte Summen und die additive Schätzung}
\label{sec:outlook-elim1}

Das gemischte $(3,4)$-Resonanzbeispiel in
\Cref{sec:hd3-approach} widerlegt eine pauschale
$q^{1+\varepsilon}$-Schätzung für jeden nichttrivialen Charakter
bei $\varepsilon<1/2$. Es widerlegt nicht jede charakterbeschränkte
Schätzung für eine bestimmte Signatur wie $(3,4,5)$.
Die schwächere Schätzung aus \Cref{lem:weil} reicht für die additive
Schranke. Verbesserte lokale Summen allein liefern nicht den
fehlenden Beyond-CRT-Transfer aus \Cref{conj:exp}.
Historische Versuche mit Selberg-Gewichten, bilinearer Dispersion,
Heath-Brown-Identitäten und Halasz-artigen Schranken
\cite{BFI1986,HB1982,HalaszPretentious} dokumentieren erfolglose
Transfers innerhalb der getesteten Formulierungen. Sie beweisen
weder eine universelle Paritätsbarriere zweiter Ordnung noch den
Ausschluss aller Standard-Siebmethoden.

\subsection{Detektion, Zählung und der kubische Fall}
\label{sec:outlook-elim2}

Primzahldetektionssätze für Polynomwerte
\cite{FI1998a,FI1998b,HB2001,HBMoroz2002,Merikoski2022}
dürfen ohne getypte Reduktion nicht als Sätze zur Zählung exakter
Lösungen importiert werden. Die Faktorisierung
$A^3+B^3=(A+B)(A^2-AB+B^2)$ verhindert die direkte Anwendung
einer Voraussetzung, die eine irreduzible kubische Form fordert;
sie beweist keine kategorische Unverträglichkeit von Detektion
und Zählung.
Der positive primitive Fall $(3,3,3)$ wird durch den klassischen
Fermat-Satz für Exponent drei ausgeschlossen. Weil seine Anzahl
null ist, gilt jede obere Schranke mit nichtnegativer rechter Seite
automatisch; die Aussage, Abstieg könne keine solche Schranke
liefern, ist falsch.
Keine ungebundene Variablenzahlschwelle wird als allgemeiner
Unmöglichkeitssatz verwendet.

\subsection{Signaturklassifikation und ihre Grenzen}
\label{sec:outlook-phase}

Das Vorzeichen von $\chi=1/x+1/y+1/z-1$ klassifiziert
Orbifold-Signaturen. Im Bereich $x,y,z\ge3$ ist $(3,3,3)$
der einzige Nullfall; die übrigen Fälle sind hyperbolisch.
Außerhalb dieses Bereichs ist $(2,2,z)$ wegen $\chi=1/z>0$
sphärisch. Das Vorzeichen allein beweist weder parametrische
primitive Familien noch ein exaktes Geschlecht einer glatten
Überlagerung, Siebanwendbarkeit oder Leerheit.
Darmon--Granville liefert Endlichkeit für jede feste hyperbolische
Signatur. Selmers Beispiele diagonaler kubischer Kurven
\cite{Selmer1951} warnen vor lokal-globalem Transfer, sind aber
kein Satz über die Fermat-Kubik.

\subsection{Eine getypte Beyond-CRT-Forschungsfrage}
\label{sec:outlook-rqprime}

Lässt sich die Produktinformation aus \Cref{lem:crt-window}
durch einen bewiesenen Korrelationssatz auf eine genau bestimmte
Zählung in einer kurzen Box übertragen?
Eine mögliche Frey-, Selmer-, Determinanten- oder geometrische
Siebroute \cite{BS2015,Bhargava2014,HeathBrown2002,KoymansSmith2024}
muss Objekte, Abbildung, Höhenvergleich, Multiplizitätsschranke,
Familieneinschränkungen und uniforme Konstanten benennen.
Ein automatischer Transfer dieser zitierten Ansätze auf allgemeine
gemischte Beal-Signaturen ist hier nicht etabliert.
Insbesondere sind weder eine universelle $H^{2/3}$-Obergrenze
noch eine hinreichende Sato--Tate-/Lang--Trotter-Implikation bewiesen.
Für die kumulative exakte Zählung ist das exponentielle Ziel
bereits Beal-stark; für eine lokale Passzählung ist es eine
andere Vermutung.

Die neuesten internen Root-Coherence-Notizen schärfen diesen empirischen Schritt zu einem Rough-Modulus-Root-Coherence-Test. Die Diagnostik zählt nicht nur lokale Pass/Fail-Ereignisse; sie trennt kleine Prim-Obstruktionen, Rough-Modulus-Dichte der Wurzelmengen $R_q(A,B)$, Unit-/Bewertungs-/Branch-Kohärenz und den verbleibenden nichtkohärenten Rest. Das ist ein experimentelles Ledger, kein Theorem. Sein Zweck ist sichtbar zu machen, ob eine globale Wurzelsektion kleiner Höhe an gewöhnlichen lokalen Obstruktionen, an fehlender Rough-Modulus-Dichte oder an einem echten nichtlokalen Kohärenzdefekt scheitert.

Ein begrenzter interner Pilot dieses Ledgers vom 22. Mai 2026 mit 250 Near-Miss-Zeilen und $Q\in\{251,503,1009\}$ fand bei allen echten Near-Misses lokale Obstruktion vor Drift: In jedem der drei Primblöcke hatten $250/250$ echte Near-Misses ein leeres lokales Wurzelset, und $0/250$ niedrige Wurzelsektionen überlebten. Die Juni-2026-Nachzugsaudits kehrten dieses negative Muster nicht um. Das S-Unit-/Frey-Shadow-Audit fand unter den fünf natürlichen $Q=61$-Survivors $0$ Q97-plus-Survivors und $0$ direkte source-bound S-Unit-/Frey-Anker. Das Q97-plus-Synthetic-Control-Audit zeigte, dass lokale Q97-plus- oder $Q=127$-Passaggregate durch registrierte Sektionen oder permutierte Branch-Maps synthetisch erzeugt werden können; diese Kontrollen brechen jedoch am Branch-Binding oder an der niedrigen globalen Sektion. Das schließt HD3' nicht. Es verschiebt die nächste empirische Aufgabe: Ein positiver Survivor muss alle drei Gates zugleich erfüllen -- Q97-plus lokales Überleben, tatsächliches Branch-Binding aus $A^x+B^y$ und eine globale Wurzelsektion kleiner Höhe oder source-bound S-Unit-/Frey-Daten -- bevor er mehr als ein diagnostischer Stresstest ist.

\begin{remark}[Status der Forschungsfrage]
Diese historischen Tests liefern ein begrenztes Diagnose-Ledger.
Sie zertifizieren weder HD3' noch den Ausschluss sämtlicher
alternativer Methoden. Der korrigierte additive Satz behält genau
den Definitionsbereich und die Großhöhenbedingungen aus
\Cref{thm:sieve}.
\end{remark}

\section{Fazit und ehrlicher Status}

HD1 und HD-PC sind elementare Normalisierungen. Das formale
HD2'-Budget bei festem $K$ folgt asymptotisch aus den angegebenen
klassischen lokalen Inputs und zertifiziert weder Wurzelidentität
noch globalen Ausschluss. HD3' bleibt offen und benötigt einen
getypten Wurzel-/Höhen-/Primzahlvertrag.
Der lokale Dichtebeweis, das additive Sieb für positive Zeilen und
die Schätzung für feste Progressionen liefern \Cref{thm:sieve}
nur in seinem angegebenen Großhöhenbereich.
Kein explizites Konstantenpaket und keine zertifizierte numerische
Schwelle werden geliefert.
Die exponentielle Schranke für die kumulative positive primitive
exakte Zählung würde die Anzahl identisch auf null setzen; ihre
Formulierung für alle $K$ ist bereits Beal-stark.
Historische numerische Tabellen, angepasste Scores, fest gesetzte
Gates und übereinstimmende Quellendatei-Hashes liefern keinen
zusätzlichen Ausschlusssatz.
Diese Korrekturfassung stuft das historische Claim-Register nicht
hoch, zertifiziert keinen externen Computerbeweis und etabliert
keinen allgemeinen Unmöglichkeitssatz für alternative Methoden.

\section{Korrekturvermerk zur Inventur für Version 1.1}
\label{sec:inventory-corrections}

Dieser Abschnitt dokumentiert den Umfang der R25-Inventurkorrekturen.
Er trennt elementare Identitäten, historische Diagnostik,
quellenberichtete Literatur und weiterhin offene arithmetische Transfers.
Originale Producer-Dateien, Ergebnistabellen, Claim-Status und
historische Prüfsiegel werden durch diese Manuskriptfassung nicht
umgeschrieben.

\subsection{Elementare Höhenidentitäten und Diagnose-Gates}
Setze für eine exakte positive Summe $U+V=H$ den Wert $u=U/H$ und
\[
 \Xi=\log\frac{H^2}{UV}=-\log[u(1-u)]\ge\log4,\qquad
 \Delta=\Xi-\log2\ge\log2.
\]
Diese Ungleichungen sind elementar und gelten auch für nichtprimitive
exakte Tupel. Feste ungleiche Anteile lassen $\Xi$ nicht divergieren;
Divergenz benötigt $\min(U,V)/H\to0$.
Die exakte Familie
$(A,B,C;x,y,z)=(3t^5,6t^5,3t^3;3,3,5)$, $t\ge1$,
hat $u=1/9$ und konstantes $\Xi$ und ist nichtprimitiv.
Sie ist kein Beal-Gegenbeispiel.
Ein IUT-/Max-Hull- oder gcd-Ausschluss benötigt eine eigene,
derzeit fehlende arithmetische Brücke; die elementare
$\Delta$-Ungleichung liefert sie nicht.
Ein als $\log H\,|\chi|$ definierter Term wächst aufgrund seiner
Definition, nicht aufgrund einer nachgewiesenen arithmetischen
Obstruktion.

Die August-/September-Diagnoseformel
$S_{\rm PH}=-\sum_q (1/q)/(1+\phi_q^2)$ ist für jede nichtleere
Primzahlmenge und jeden reellen Input $\phi_q$ negativ.
Ihr Vorzeichen kann ein primitives Gegenbeispiel nicht unterscheiden.
Das historische Branch-Binding-Gate ist ein konstant wahrer Wert,
keine implementierte Bindungsprüfung.
Das historische Survivor-Prädikat verwendet Primitivität,
Exponentenbereich und Exaktheit; es verwendet die drei Diagnose-Gates
nicht. Sein Survivor-Zählwert null validiert daher den vorgeschlagenen
Ausschlussmechanismus nicht.

Der historische Replay über 25 Zeilen enthält elf primitive Zeilen
im Exponentenbereich: zwei parabolische $(3,3,3)$-Zeilen und neun
hyperbolische Zeilen. Die benannte Gruppe von zehn hyperbolisch-
primitiven Zeilen enthält eine parabolische Zeile.
Eine der fünf verankerten Kontrollen hat Exponent zwei und ist
nicht exakt. Somit haben $0/11$, $0/10$ und $4/5$ verschiedene
Nenner und Geltungsbereiche; sie sind keine austauschbaren
Erfolgszertifikate. Verschachtelte Primzahlschranken verwenden
Beobachtungen erneut. Ein $\eta$-Variationskoeffizient ist
beschreibend, kein repliziertes Experiment, Konfidenzintervall
oder Signifikanztest.

\subsection{Die Juli-Funktion und die Provenienzgrenze}
Die getrennte Juli-Diagnostik hat die Form
\[
 F(w)=\sum_{q,r}\frac{a_{q,r}}{w-e^{2\pi i r/q}},
 \qquad a_{q,r}>0,
\]
über nichtleere Wurzelmengen. Eine Auswertung außerhalb der Pole
bei $w=0.5+0.5i$ beweist die Herglotz-Eigenschaft nicht;
diese verlangt Holomorphie in der oberen Halbebene.
Bei der exakten nichtprimitiven Kontrolle $2^3+2^3=2^4$ enthalten
die Wurzeln modulo fünf $r=1,2,3,4$, sodass Pole in der oberen
Halbebene mit positiven, nicht aufgehobenen Residuen entstehen.
Das widerlegt diese Herglotz-Behauptung; es überträgt das getrennte
August-Vorzeichenargument nicht auf Juli.

Die Inventur fand fünf PDF-Hashes, die zum JSON-Validierungsmanifest
und den geprüften lokalen Dateien passen, aber widersprechende
Hashfelder im begleitenden Markdown-Validierungsvermerk.
Ein lokaler Byte-Identitätscheck ist weder ein externes
Provenienzaudit noch eine Beweisprüfung. Historische Dateien
und Siegel bleiben erhalten; dieser Korrekturtext versiegelt
sie nicht stillschweigend neu.

\subsection{Begrenzter Literatur-Nachtrag}
Chocians Preprint \cite{Chocian2026} berichtet einen
computerunterstützten Ausschluss nichtverschwindender teilerfremder
ganzzahliger Lösungen der Signatur $(3,5,7)$.
Seine arXiv-Identität ist bestätigt. Beweis und Replay-Begleitpaket
wurden hier nicht unabhängig validiert; sie zertifizieren nicht
die allgemeine Verifier-Architektur dieses Papers.

Sahoo \cite[Satz~1.10, Proposition~1.12]{Sahoo2026} setzt
Vermutungen~2.4 und~2.5 voraus: Bei quadratfreiem $d\ge5$,
$-d\not\equiv1\pmod8$, Koeffizienten $\{\pm2^r\}$ und der
Quellenbedingung
$\alpha_{\mathfrak p}\beta_{\mathfrak p}\gamma_{\mathfrak p}\ne0$
schließt er hinreichend große \emph{gemeinsame Primexponenten}
in $W_{L,\mathfrak p}$ aus, wobei
$L=\mathbb Q(\sqrt{-d})$, $\mathfrak p\mid2$ und
$\mathfrak p\mid abc$.
Die Dichte $5/6$ ist relativ zu quadratfreien $d$.
Dies ist quellenberichteter Umfang, kein allgemeiner HD3'-Satz
für gemischte Exponenten oder eine uniforme Turm-/Koerzivitätsabschätzung.

Pasten \cite[Korollare~1.2--1.3]{Pasten2026} berichtet
$h(E)\ll N\log\log N$ für elliptische Kurven über $\mathbb Q$
und $h(E)\ll_S N$ für Kurven, die außerhalb eines festen $S$
semistabil sind. Die Aussagen betreffen Faltings-Höhe und
Konduktor im Sinne der Quelle.
Sie beweisen weder die Szpiro-Vermutung noch entfernen sie
automatisch Modulformen in einem Beal-Transfer.
Die benötigten uniformen Höhen- und S-Unit-Vergleiche bei
variierenden Primzahlmengen oder Körpern bleiben offen.

\appendix

\section{Empirische Daten und Reproduzierbarkeit}
\label{app:data}

Der historische Near-Miss-Korpus mit 5000 Zeilen wurde mit
\path{evidence/scan_near_power_fractures.py}, den Parametern
\texttt{--max-base 250}, \texttt{--min-exp 3}, \texttt{--max-exp 7} und
\texttt{--keep 5000}, \texttt{--max-rel-gap 1e-4} erzeugt. Das variable-$Q$-Experiment
verwendet die ersten 1000 nach relativer Lücke geordneten Zeilen und
\path{evidence/run_variable_Q_sieve.py} für
$Q\in\{251,503,1009,2003\}$. Die dokumentierte Laufzeit von etwa
14 Minuten auf einem Server mit zwei Kernen ist historisch und kein
neuer Benchmark.

Das kuratierte Beal-Quellen- und Datenpaket ist unter
\url{https://github.com/research-line/functional-stability-theory/tree/codex/beal-inventurernte-v1.1/fst-mathematics/beal-height-dominance}
verfügbar. Es enthält beide Manuskriptquellen, die Disclosure-Eingaben,
acht historische Diagnoseskripte und 17 unveränderte CSV-Dateien.
Die README des Pakets dokumentiert Datensatzauswahl, Kontrollklassen,
Wiederholungsbefehle und Claim-Grenzen. Das Manifest bindet die
bereitgestellten Artefakte über SHA-256; die zugehörigen Tests prüfen
Artefaktintegrität und die berichteten endlichen Diagnoseaggregate,
keinen Beweis der Bealschen Vermutung.

Der Root-Branch-Drift-Pilot vom 22.~Mai liefert drei Ledger mit jeweils
750 Zeilen: 250 echte Near-Misses, 250 zufällige Root-Kontrollen und
250 Exact-Power-Kontrollen für jedes $Q\in\{251,503,1009\}$.
Der Nachzug vom 1.~Juni verwendet $Q\in\{7,13,31,61\}$ für dieselbe
Auswahl von 250 Zeilen und dieselben drei Kontrollklassen.
Der rank-8-Mikroaudit vom 2.~Juni identifiziert den ersten leeren
Kubikroot-Kanal bei $q=37$; die Above-threshold-Suche im Korpus mit
5000 Zeilen hat fünf lokale Passes bei $Q=61$ und keine bei $Q=97$.
Die S-Unit-Shadow-Tabelle vom 22.~Juni und die Tabelle synthetischer
Kontrollen vom 23.~Juni sind als historische Klassifikationen enthalten.
Synthetische registrierte, verschobene oder geschuffelte Sektionen
sind Kontrollen und können keine Branch-Bindung an die ursprüngliche
arithmetische Gleichung begründen.
Diese Dateien tragen die deskriptive Prüfspur aus
\cref{tab:root-diagnostics}; sie begründen weder HD3', Root-Kohärenz,
statistische Signifikanz noch eine effektive Höhenschwelle.
Interne Beweis- und Reviewnotizen sowie fremde Quellen-PDFs sind
nicht Teil des öffentlichen Pakets.

% =============================================================================
% AI DISCLOSURE -- PIPELINE-STANDARD, DEUTSCHE FASSUNG (seit 2026-08-17, T-20260817-02)
% =============================================================================
% Sprachreine Fassung fuer rein deutschsprachige Manuskripte. Text 1:1 aus dem
% deutschen Block von ai_disclosure_STANDARD.tex uebernommen (keine inhaltliche
% Aenderung) -- nur die Ueberschrift ist jetzt sprachrein und der englische
% Block entfaellt. Die bilinguale Datei ai_disclosure_STANDARD.tex bleibt
% unveraendert bestehen (Bestandsverweise duerfen weiter darauf zeigen).
%
% ABWEICHEN NACH UNTEN: Wurde in einem Paper tatsaechlich deutlich weniger
% KI eingesetzt, wird NICHT auf eine andere Datei gewechselt, sondern die
% Formulierung hier abgeschwaecht -- z. B. "in erheblichem Masse in allen
% Phasen" ersetzen durch die tatsaechlich zutreffenden Phasen, und aus der
% Qualitaetssicherungs-Aufzaehlung streichen, was nicht stattfand.
% Der Grundsatz bleibt: ehrlich benennen, was war -- nicht abgrenzen.
%
% Verwendung (nur fuer deutschsprachige Manuskripte):
%   \input{../../_templates/disclosure/ai_disclosure_STANDARD_de}
% =============================================================================

\section*{KI-Offenlegung}

\noindent
KI-Systeme wurden in erheblichem Ma\ss{}e in allen Phasen dieser Arbeit
eingesetzt: Konzeption und Forschungsdiskussion, Literatur- und
Quellenarbeit, mathematische Analyse und Beweisarbeit, Erstellung und
Pr\"ufung der Rechenskripte, Ausf\"uhrung und Auswertung der
Berechnungsl\"aufe, Text, \"Ubersetzung und Satz.

\medskip

\noindent
Qualit\"atssicherung: Einsatz mehrerer unabh\"angiger Modelle
verschiedener Anbieter mit Gegenpr\"ufungen und adversarialen KI-Reviews;
nachvollziehbare, versionierte Rechenskripte mit gespeicherten
Ergebnisdateien und Pr\"ufsummen; fortlaufende Beweis- und
Registerdokumente; strikte Claim-Hygiene (Behauptungen nur mit
dokumentiertem Beleg, Nachtr\"age statt stiller Korrekturen); Changelogs
\"uber alle Versionen.


\begingroup
\small
\setlength{\itemsep}{0pt}
\setlength{\parskip}{0pt}
\setlength{\parsep}{0pt}
\begin{thebibliography}{99}

\bibitem{BFI1986}
E.~Bombieri, J.~B.~Friedlander, and H.~Iwaniec, ``Primes in arithmetic progressions to large moduli,''
\emph{Acta Mathematica} 156 (1986), 203--251. doi:
\href{https://doi.org/10.1007/BF02399204}{10.1007/BF02399204}.

\bibitem{Bhargava2014}
M.~Bhargava, ``The geometric sieve and the density of squarefree values of invariant polynomials,''
arXiv:1402.0031 (2014). \url{https://arxiv.org/abs/1402.0031}.

\bibitem{Bombieri1971}
E.~Bombieri, ``A note on the large sieve,''
\emph{Acta Arithmetica} 18 (1971), no.\ 1, 401--404. doi:
\href{https://doi.org/10.4064/aa-18-1-401-404}{10.4064/aa-18-1-401-404}.

\bibitem{BS2015}
M.~Bhargava and A.~Shankar, ``Binary quartic forms having bounded invariants, and the boundedness of the average rank of elliptic curves,'' \emph{Annals of Mathematics} 181 (2015), 191--242. arXiv:1006.1002.
doi: \href{https://doi.org/10.4007/annals.2015.181.1.3}{10.4007/annals.2015.181.1.3}.

\bibitem{DarmonGranville1995}
H.~Darmon and A.~Granville, ``On the equations $z^m=F(x,y)$ and $Ax^p+By^q=Cz^r$,''
\emph{Bulletin of the London Mathematical Society} 27 (1995),
no.\ 6, 513--543. doi:
\href{https://doi.org/10.1112/blms/27.6.513}{10.1112/blms/27.6.513}.

\bibitem{FI1998a}
J.~B.~Friedlander and H.~Iwaniec, ``The polynomial $X^2+Y^4$ captures its primes,''
\emph{Annals of Mathematics} 148 (1998), 945--1040.
doi: \href{https://doi.org/10.2307/121034}{10.2307/121034};
arXiv:math/9811185.

\bibitem{FI1998b}
J.~B.~Friedlander and H.~Iwaniec, ``Asymptotic sieve for primes,''
\emph{Annals of Mathematics} 148 (1998), 1041--1065.
doi: \href{https://doi.org/10.2307/121035}{10.2307/121035};
arXiv:math/9811186.

\bibitem{FouvryKatz2001}
É.~Fouvry and N.~M.~Katz, ``A general stratification theorem for exponential sums, and applications,''
\emph{Journal für die reine und angewandte Mathematik (Crelle)} 540 (2001), 115--166. doi:
\href{https://doi.org/10.1515/crll.2001.082}{10.1515/crll.2001.082}.

\bibitem{HalaszPretentious}
G.~Halász, \glqq Über die Mittelwerte multiplikativer zahlentheoretischer Funktionen\grqq{},
\emph{Acta Mathematica Academiae Scientiarum Hungaricae} 19 (1968), no.\ 3--4, 365--403. doi:
\href{https://doi.org/10.1007/BF01894515}{10.1007/BF01894515}.

\bibitem{HB1982}
D.~R.~Heath-Brown, ``Prime numbers in short intervals and a generalized Vaughan identity,''
\emph{Canadian Journal of Mathematics} 34 (1982), no.\ 6, 1365--1377.
doi: \href{https://doi.org/10.4153/CJM-1982-095-9}{10.4153/CJM-1982-095-9}.

\bibitem{HB2001}
D.~R.~Heath-Brown, ``Primes represented by $x^3+2y^3$,''
\emph{Acta Mathematica} 186 (2001), no.\ 1, 1--84. doi:
\href{https://doi.org/10.1007/BF02392715}{10.1007/BF02392715}.

\bibitem{HBMoroz2002}
D.~R.~Heath-Brown and B.~Z.~Moroz, ``Primes represented by binary cubic forms,''
\emph{Proceedings of the London Mathematical Society} 84 (2002), no.\ 2, 257--288. doi:
\href{https://doi.org/10.1112/plms/84.2.257}{10.1112/plms/84.2.257}.

\bibitem{HeathBrown2002}
D.~R.~Heath-Brown, ``The density of rational points on curves and surfaces'' (with an appendix by J.-L.~Colliot-Thélène),
\emph{Annals of Mathematics} 155 (2002), no.\ 2, 553--598. doi:
\href{https://doi.org/10.2307/3062125}{10.2307/3062125}.

\bibitem{IrelandRosen}
K.~Ireland and M.~Rosen, \emph{A Classical Introduction to Modern Number Theory}, 2nd ed., Graduate Texts in Mathematics 84, Springer, 1990.
doi: \href{https://doi.org/10.1007/978-1-4757-2103-4}{10.1007/978-1-4757-2103-4}.

\bibitem{IwaniecKowalski2004}
H.~Iwaniec and E.~Kowalski, \emph{Analytic Number Theory},
American Mathematical Society Colloquium Publications 53, 2004. doi:
\href{https://doi.org/10.1090/coll/053}{10.1090/coll/053}.

\bibitem{KatzTwistedL}
N.~M.~Katz, \emph{Twisted $L$-Functions and Monodromy},
Annals of Mathematics Studies 150, Princeton University Press, 2002.

\bibitem{KoymansSmith2024}
P.~Koymans and A.~Smith, ``Sums of rational cubes and the $3$-Selmer group,'' arXiv:2405.09311 (2024).
\url{https://arxiv.org/abs/2405.09311}.

\bibitem{Mauldin1997}
R.~D.~Mauldin, ``A Generalization of Fermat's Last Theorem: The Beal Conjecture and Prize Problem,''
\emph{Notices of the AMS} 44 (1997), no.\ 11, 1436--1437.
\url{https://www.ams.org/notices/199711/beal.pdf}.

\bibitem{Merikoski2022}
J.~Merikoski, ``A cubic analogue of the Friedlander--Iwaniec spin over primes,''
\emph{Mathematische Zeitschrift} 300 (2022), 2809--2835.
doi: \href{https://doi.org/10.1007/s00209-021-02878-5}{10.1007/s00209-021-02878-5};
arXiv:2012.05675.

\bibitem{Selmer1951}
E.~S.~Selmer, ``The Diophantine equation $ax^3+by^3+cz^3=0$,''
\emph{Acta Mathematica} 85 (1951), 203--362. doi:
\href{https://doi.org/10.1007/BF02395746}{10.1007/BF02395746}.

\bibitem{SiksekStoll2012}
S.~Siksek and M.~Stoll, ``Partial descent on hyperelliptic curves and the generalized Fermat equation $x^3+y^4+z^5=0$,''
\emph{Bulletin of the London Mathematical Society} 44 (2012), no.\ 1, 151--166.
doi: \href{https://doi.org/10.1112/blms/bdr086}{10.1112/blms/bdr086};
arXiv:1103.1979.

\bibitem{Weil1948}
A.~Weil, \emph{Sur les courbes algébriques et les variétés qui s'en déduisent},
Actualités scientifiques et industrielles 1041, Hermann, Paris, 1948.

\bibitem{Xu2020}
J.~Xu, ``Stratification for multiplicative character sums,''
\emph{International Mathematics Research Notices} 2020, no.\ 10, 2881--2917.
doi: \href{https://doi.org/10.1093/imrn/rny096}{10.1093/imrn/rny096};
arXiv:1709.01663.

\bibitem{KedlayaANT}
K.~S.~Kedlaya, \emph{Notes on Analytic Number Theory},
online author notes, Theorems~18.1 and~4.12, accessed October~1, 2026.
\url{https://kskedlaya.org/ant/chapter-18.html};
\url{https://kskedlaya.org/ant/chap-primes-in-ap.html}.

\bibitem{Chocian2026}
P.~Chocian, \emph{The Primitive Generalized Fermat Equation
$x^3+y^5=z^7$: A computer-assisted proof},
arXiv:2609.26996 (2026), source-reported preprint.
\url{https://arxiv.org/abs/2609.26996}.

\bibitem{Pasten2026}
H.~Pasten, \emph{Improved bounds for Szpiro's conjecture},
arXiv:2609.17390 (2026), source-reported preprint.
\url{https://arxiv.org/abs/2609.17390}.

\bibitem{Sahoo2026}
S.~Sahoo, \emph{Generalized Fermat equation over number fields},
arXiv:2608.28117 (2026), source-reported preprint.
\url{https://arxiv.org/abs/2608.28117}.

\end{thebibliography}
\endgroup

\end{document}
